\newpage

% ============================================
% PÁGINA DE ABSTRACT (sem numeração)
% ============================================
\thispagestyle{empty}

\begin{center}
    \large\textbf{ABSTRACT}
\end{center}

\vspace{0.5cm}

\noindent Compilers are a fundamental subject in Computer Science curricula, yet parsing algorithms are notoriously difficult to teach and learn due to their abstract nature and multiple intermediate steps. This work presents VANSI (Visualization of Syntax Analysis Algorithms), an interactive, cross-platform educational tool designed to support the teaching of LL(1), SLR, and LR(1) parsing algorithms. Developed using the Svelte framework and packaged for desktop and mobile via Tauri and Capacitor, VANSI provides step-by-step algorithm execution with animations, interactive automata, parse trees, parsing tables, and pseudocode with breakpoints. The tool was evaluated through a mixed-methods study involving 23 students from the Federal University of Ceará, Quixadá Campus, combining automated interaction logs with structured questionnaires. Results indicate high usability scores (mean 4.2/5) and strong pedagogical effectiveness (mean 4.4/5), with 91\% of participants expressing satisfaction and 90\% reporting intention to reuse the tool. VANSI addresses gaps identified in existing tools by offering greater interactivity and detailed step-by-step visualization of internal algorithmic processes.

\vspace{0.5cm}

\noindent\textbf{Keywords:} compiler construction. parsing algorithms. algorithm visualization. educational technology. Svelte.
